\section{Practices for Each Role}
\label{sec:rolepractices}

The fourteen practices in this section are specific to roles. Each adapts the
tradition named in Table~\ref{tab:roles} to the conditions of HRI.

\subsection{R1 Perceiver}

\begin{practice}{bp:perceiver}{Evaluate a perceiving model as a measurement instrument}
Report per-class performance against a reference produced by more than one
qualified rater, express model--human agreement relative to human--human
agreement, and, where the output will inform decisions about a person, meet
the standard expected of an assessment tool.
\end{practice}

Language and vision-language models are used to infer what a robot cannot
sense directly: emotion from a face or a transcript, engagement from video,
intent from an ambiguous request, wellbeing from a child's narrative. This is
classification or scoring, and the intuition that evaluating it looks like
evaluating a classification model is right, with three qualifications that
the affective-computing and psychometric literatures supply.

First, the reference labels are themselves judgements. For affect,
engagement, and rapport there is no instrument-independent truth, only what
trained observers or the person themselves report, and these disagree. The
model cannot be expected to agree with the reference more than the reference
agrees with itself, so human--human agreement must be measured and reported
as the ceiling (\bpr{bp:human}). A reference produced by a single rater, a
frequent economy, makes the ceiling unknown and the model's score
uninterpretable; the authors of the child wellbeing study name their single
psychologist rater as a principal limitation for this reason
\cite{abbasi2025robotled}.

Second, class imbalance is the norm, so accuracy is uninformative and
per-class precision, recall, and agreement are required (\bpr{bp:human},
\bpr{bp:subgroup}).

Third, what the model is given matters as much as which model it is. A model
may receive raw video, a transcript, or features extracted by an upstream
component, such as valence and arousal values computed from facial
expressions and passed to a language model as numbers \cite{mehra2025how}.
Each arrangement has a different error profile, and the evaluation should be
of the arrangement used on the robot, with the robot's sensors, viewpoint,
and lighting, and not of the model on a curated benchmark. The gap can be
large: multimodal models rating emotional arousal from video approached
human-level reliability on laboratory recordings but correlated at best
moderately with human ratings of real-world footage, and showed
gender-differential bias \cite{licht2025computational}.

When a Perceiver's output is used to draw conclusions about an individual
(screening, assessment, adaptation of an intervention), the relevant standard
is that of psychological testing: evidence of reliability, of validity
against an external criterion, and of fairness across groups, for the
intended use and population \cite{aera2014standards}. A model that performs
``reasonably'' on average does not meet it.

\subsection{R2 Representer}

\begin{practice}{bp:embed}{Show that structure derived from embeddings is not an artefact of the embedding}
When embeddings are used to cluster, compare, or measure similarity, repeat
the analysis with more than one embedding model, calibrate similarity against
an empirical null, test the stability of clusters, and have people check the
resulting structure.
\end{practice}

Embeddings have become a standard analytic instrument in HRI as elsewhere:
utterances are embedded, clustered into themes, and the themes labelled by a
language model; or the similarity of two sets of texts is taken as a measure
of how alike two conditions, partners, or populations are. The approach
scales to corpora that could not be coded by hand, and it has produced
findings about what people disclose to robots that would otherwise be
unavailable \cite{chiang2025selfdisclosure,laban2025what}. Its validity rests
on four checks.

\emph{Triangulation.} Cosine similarity is defined relative to an embedding
space. Contextual embedding spaces are anisotropic, so that even unrelated
texts have high cosine similarity in some models
\cite{ethayarajh2019how}, whereas models trained with a sentence-similarity
objective spread texts more widely \cite{reimers2019sentence}. The same pairs
of disclosures were found to have mean similarity below .50 in one family and
above .80 in two others \cite{chiang2025selfdisclosure}. Conclusions should
therefore be reported for at least two embedding families, and those that
hold in only one should be flagged.

\emph{Null calibration.} For the same reason, fixed thresholds are not
portable. The reference point for an observed similarity is the distribution
of similarity between randomly paired items in the same space, or under
permutation of the grouping variable, and the result is best expressed as a
standardised difference from that null.

\emph{Stability.} A clustering is one of many that an algorithm could have
returned. The number of clusters, the random initialisation, and the sample
all move the partition. Reporting how well the clusters are recovered across
seeds and bootstrap resamples, for example by the mean adjusted Rand index or
cluster-wise Jaccard similarity \cite{hennig2007cluster}, distinguishes
structure in the data from structure in the run.

\emph{Human check.} Automatically generated cluster labels are fluent whether
or not the cluster is coherent. Topic-model research met this problem and
devised intrusion tasks, in which people try to spot the item that does not
belong \cite{chang2009reading}. The same test applies to embedding clusters,
alongside a conventional check in which two coders assign a sample of items
to the labelled themes. An automated check that each generated description
lies closest to its own centroid, as used in the studies above, is a useful
necessary condition and costs nothing; it is not a substitute.

\begin{practice}{bp:memory}{Evaluate retrieval and memory as retrieval}
When embeddings select what the robot remembers or retrieves, report
retrieval quality and the rate of retrieving material that should not have
been retrieved, including material about a different person.
\end{practice}

Long-term interaction increasingly relies on retrieval: past conversations or
user facts are embedded, and the most similar are inserted into the prompt.
The retrieval step is an information-retrieval system and can be evaluated as
one, with precision and recall against human judgements of relevance, and
with the faithfulness of the generated response to what was retrieved
\cite{es2024ragas}. HRI adds a failure with social and legal weight:
retrieving one person's disclosure into a conversation with another, in a
household or care home where a robot serves several people. The rate of such
cross-person retrieval should be measured directly.

\subsection{R3 Decider}

\begin{practice}{bp:decider}{Evaluate decisions against expert decisions, and evaluate whether the model knows when not to decide}
Compare the model's choices with those of qualified people in the same
situations using agreement statistics, report violations of stated
constraints, and measure whether the model's abstentions and requests for
help are calibrated to its uncertainty.
\end{practice}

Models decide which skill to execute, which tool to call, and what to do
next in a task plan \cite{ahn2022do,liang2023code}; in social robots they
also decide when to take a turn \cite{skantze2025applying}, and whether a
situation calls for help, silence, or intervention
\cite{tanneberg2024help}. Such decisions are sometimes verifiable (a plan is
executable or not) and sometimes normative (it was or was not appropriate to
interrupt). For the normative kind, the evaluation is a comparison with
people's decisions in the same situations, subject to everything said about
human references above. GRACE shows what such an evaluation looks like when
the reference is taken seriously \cite{dogan2025grace}. It predicts how
socially appropriate each of several household actions (vacuuming, carrying
warm food, starting a conversation) would be in a given scene, and it
evaluates five language models, open-weight and proprietary, against
appropriateness scores given by thousands of human annotators, using error
and two correlation measures, one of which, the concordance correlation
coefficient, penalises systematic offset (Table~\ref{tab:irr}). Two features
make it an exemplar for this role. It does not treat the human scores as a
single truth: scenes are first separated into those on which annotators
agree and those on which they do not, and model performance is reported for
each, which showed that the models tracked human judgement more closely
where people agreed with one another and least closely where they did not.
And it uses the human explanations that accompany the scores to correct the
language model's predictions in the uncertain scenes, which improved
agreement with human scores over the language model alone (for the
strongest model, concordance rose from .23 to .48). The lesson for
evaluation is general: where a Decider's task is normative, performance
should be reported conditional on human agreement, because a model cannot be
right about a scene on which people are divided, only more or less
representative of them (Section~\ref{sec:perspectivist}).

Two cautions apply. A model that predicts human preferences well on
vignettes has been proposed as a zero-shot human model for planning
\cite{zhang2023large}, but apparent social reasoning can be brittle:
performance on theory-of-mind style judgements in HRI scenarios degraded
under perturbations that should not have mattered \cite{verma2024theory}.
Decisions should therefore be tested under perturbation of irrelevant
context as well as on canonical cases. And a Decider should be evaluated on
its uncertainty, because the safe response to ambiguity in HRI is usually to
ask. Conformal prediction has been used to make an LLM planner request human
help at a rate tied to a specified level of task success
\cite{ren2023robots}. Whether or not such machinery is used, the rate at
which the robot asked when it should have, and did not ask when it need not
have, is a reportable property with a direct bearing on the interaction.

\subsection{R4 Speaker}

\begin{practice}{bp:rubric}{Rate delivered language on a rubric that covers what embodiment adds}
Have trained raters rate the robot's utterances for relevance, correctness,
safety, and transparency, and also for grounding in the robot's body and
situation, consistency of persona, fitness for the population, and fitness
for speech.
\end{practice}

The four criteria named in the discussion that prompted this article
(relevance, correctness, safety, transparency) are the core of any rubric
for user-facing language, and each needs an operational definition, since
reviews of NLG evaluation show that criteria with the same name are defined
in incompatible ways across papers \cite{howcroft2020twenty}. For a robot
the rubric needs four further criteria that text-only evaluation omits.

\emph{Embodied grounding.} Does the utterance cohere with what the robot can
perceive and do? LLM-driven robots can claim capabilities they lack
\cite{ashok2026why}, and they hallucinate and offer outdated information in
conversation \cite{irfan2025between}. This is hallucination in the
NLP sense \cite{ji2023survey}, but its reference is the robot's own body and
history, which no general-purpose factuality metric checks.

\emph{Persona consistency.} Does the robot remain the character it was
designed to be, across turns and sessions? Drift of persona is a known
property of long conversations with language models and was among the main
concerns of HRI experts consulted on prompt design \cite{ashok2026why}.

\emph{Fitness for the population.} Is the language appropriate for the
person addressed: a child, an older adult, someone disclosing distress? For
specialist content the raters must be specialists, and psychosocial risk
(manipulation, psychological harm, privacy violation, discrimination, insult)
deserves its own dimensions \cite{luo2026dialogguard}.

\emph{Fitness for speech.} Is the utterance of a length, structure, and
register that works when spoken by this voice at this speed? A response
that reads well can be unlistenable.

\begin{practice}{bp:humanref}{Compare with what a qualified human would have said}
Where a claim concerns the quality of the robot's language, collect responses
to the same contexts from qualified people and have raters blind to source
rate both.
\end{practice}

Absolute ratings are hard to interpret: is 3.9 out of 5 for empathy good? A
human reference gives the scale meaning. The strongest design presents
raters with model-generated and human-authored responses to identical
contexts, unlabelled and in random order, and reports the difference with an
interval. If the claim is that the model is as good as the human, the
analysis is an equivalence test (Section~\ref{sec:equivalence}). The human
authors should be the people whose performance is the relevant benchmark
(counsellors for a counselling robot, teachers for a tutoring robot), and
the comparison should be on the rubric's dimensions and not on whether
raters can tell which is which, since indistinguishability to lay readers is
weak evidence of quality \cite{clark2021all}. Where matched human responses
cannot be collected, comparison with an existing corpus of professional
practice is a scalable approximation, with limits discussed in
Section~\ref{sec:alignment}.

\subsection{R5 Behaviour Author}

\begin{practice}{bp:observers}{Show that generated behaviour is perceived as intended}
Test generated expressive behaviour with observers who are not told the
intended meaning, report recognition against chance with the confusion
pattern, and compare with behaviour authored by an expert.
\end{practice}

When a model generates a gesture, a facial expression, or a sequence of
motions to express an internal state \cite{mahadevan2024generative,
antony2025xpress}, the question is the one that HRI has always asked of
designed behaviour: do people read it as intended? The method is established:
naive observers view the behaviour and label it or rate it on dimensions such
as valence and arousal, and recognition is compared with chance. A recent
study of a low-degree-of-freedom robot is a reminder of what such tests tend
to find: observers often recovered the broad affective meaning of an
expression but not its specific label, and recognition varied widely by
emotion \cite{rogel2026emotional}. Generated behaviour should be held to the
same test before it is used as a manipulation, and the plausibility of the
model's textual rationale for a behaviour should not be mistaken for evidence
about the behaviour. Because generation is stochastic, the test should cover
several generations per intended meaning (\bpr{bp:consistency}).

\begin{practice}{bp:code}{Verify generated programs before they run on a robot near people}
Report the rate at which generated code or parameters were executable,
correct, and within safety constraints, and what stood between generation and
execution.
\end{practice}

Models also write robot programs, from end-user tools that turn natural
language into applications \cite{karli2024alchemist} to policies expressed as
code \cite{liang2023code}. Software engineering supplies the standard:
functional correctness on held-out tasks, and static or formal checks of
properties that must always hold. HRI has methods for verifying social
properties of interaction programs \cite{porfirio2018authoring}, and
constraint layers that reject plans violating stated safety rules have been
built for LLM-driven robots \cite{yang2024plug}. The evaluation should report
how often generated programs failed each check, not only the performance of
those that passed.

\subsection{R6 Actor}

\begin{practice}{bp:trials}{Evaluate action policies with enough trials, stated conditions, and reliable adjudication}
Report the number of trials, how initial conditions were chosen and varied,
the success criterion fixed in advance, success rates with intervals, who or
what judged success and how reliably, and failures by type.
\end{practice}

The robot-learning community has articulated best practices for evaluating
learned policies, including specifying initial conditions, running enough
trials for statistical conclusions, reporting failure modes, and separating
evaluators from developers \cite{kressgazit2024robot}, and has developed
statistical tools for bounding performance from few trials
\cite{vincent2024how} and for comparing policies with the fewest trials
needed \cite{snyder2025your}. Simulation can rank real-world manipulation
policies when it is built to correlate with reality \cite{li2024evaluating},
and distributed and automated real-world evaluation is emerging for
generalist policies \cite{atreya2025roboarena,zhou2025autoeval}. HRI studies
using VLAs should adopt these practices as given. The addition from the
psychological tradition is that success in a physical trial is a human
judgement and should be treated as one: defined in a written criterion,
adjudicated by two raters from video, with agreement reported. When a
vision-language model is used as a success detector \cite{du2023vision}, it
is a Judge and must be validated as such (\bpr{bp:judge}).

\begin{practice}{bp:people}{Put people in the evaluation of action}
Sample the variation that people introduce (their phrasing, positions,
movements, and interruptions) and measure outcomes that matter to them
(safety margins, legibility, comfort, effort) in addition to task success.
\end{practice}

Benchmarks for action models are overwhelmingly about objects on tables. A
policy that will act around people must be evaluated around people, for two
reasons. People are a source of distribution shift that object-centred
evaluation does not sample: they phrase requests in their own way, move
while the robot moves, reach in, change their minds. And task success is not
the outcome that matters to them. A handover that succeeds after a hesitant,
illegible approach, or a navigation that reaches its goal by cutting across
a conversation, is a failure of interaction that a success rate records as a
success; a study of video-language models as critics of agent behaviour was
motivated by exactly this gap between task completion and acceptable
behaviour \cite{guan2024task}. HRI has the constructs and measures needed,
among them the legibility and predictability of motion
\cite{dragan2013legibility}, perceived safety, and proxemic norms, and the
safety-engineering tradition supplies physical limits for collaborative
operation that hold regardless of what the policy intends
\cite{iso15066}. Because VLAs inherit the language interface of their
backbone, they also inherit its vulnerabilities, and instructions that would
cause unsafe actions belong in the test set (Section~\ref{sec:redteam}).

\subsection{R7 Annotator}

\begin{practice}{bp:coder}{Treat a model that codes data as a coder}
Develop the codebook with human coders, have at least two of them code a
held-out subset, report model--human agreement next to human--human agreement
with the same coefficient, and freeze the prompt before the model sees the
data on which conclusions rest.
\end{practice}

Language models can match or exceed crowd workers on some annotation tasks
\cite{gilardi2023chatgpt} and perform well on many computational social
science tasks while remaining below fine-tuned models and human experts on
others \cite{ziems2024can}. The methodological literature that followed is
consistent on what this licenses: performance varies by task and by prompt
and cannot be assumed, so every application requires validation against
human-coded data from the same corpus
\cite{pangakis2023automated,tornberg2024best}. For deductive qualitative
coding, combining an expert-developed codebook with a language model has
produced fair to substantial agreement with experts \cite{xiao2023supporting};
agreement varies by code as well as by model, with one study reporting
$\kappa\geq.79$ against a human gold standard for three of nine codes and
$\kappa\geq.60$ for eight of nine with a stronger model, against a mean of
.34 with a weaker one \cite{dunivin2024scalable}. The better reports place
the model's figure next to the human one: coding open-text survey answers
into fifty categories, the best model was correct on 93.9\% of items and a
second human coder on 94.7\% \cite{mellon2024do}; and in replications of
fourteen coding tasks, recent models agreed with attentive human coders at
rates comparable to agreement among humans, with disagreement between models
concentrated on the texts that also divided people, and with clarification
of the codebook raising both \cite{nakamura2026observational}.

The content-analytic framing makes the requirements simple. The model is an
additional coder. A coder is trained on material that is then set aside,
tested on fresh material, and admitted to the study if agreement on the fresh
material is adequate. Tuning the prompt on the same items used to report
agreement is the equivalent of training and testing a coder on the same
transcripts, and it inflates the estimate. The human--human coefficient on
the validation subset is the benchmark, and the subset must be large enough
for the coefficients to be estimated with useful precision
(Table~\ref{tab:kappa}). For multimodal annotation, where a vision-language
model codes gaze, gesture, or engagement from video, the requirements are the
same and published validations are fewer. An analysis prompted by the
finding that 125 of 1{,}702 papers at one HCI conference annotated video
found that a vision-language model alone approached human accuracy on
average across fifteen benchmark tasks, and that people verifying the
model's output did better than either alone while halving annotation time
\cite{shen2026can}; agreement on clinically meaningful judgements in HRI
has been much lower \cite{abbasi2025robotled}. Checks of this kind are
already made in HRI and are cheap to strengthen. In GRACE, for example, a
language model assigned annotators' free-text explanations to categories
that a researcher had derived from a sample of the data, and the researcher
then hand-labelled a further 200 randomly chosen explanations, which matched
the model's labels in 95\% of cases \cite{dogan2025grace}. That is a held-out
check on a stated sample, reported with its size. A second human labeller
and a chance-corrected coefficient per category would turn it into the full
form of this practice.

Not all qualitative analysis is coding against a fixed scheme. In
interpretive approaches such as reflexive thematic analysis, the analyst's
subjectivity is a resource and inter-rater statistics are a category mistake
\cite{braun2019reflecting,mcdonald2019reliability}. A model can assist such
analysis, but claims that it has ``performed'' the analysis should be
resisted; what can be reported is how the model was used, and how the
analysts checked and departed from what it offered.

\begin{practice}{bp:inference}{Carry annotation error into the statistics}
When model-produced labels enter a statistical analysis, use a method that
accounts for their error using the human-labelled subset, or show that
conclusions are robust to it.
\end{practice}

Even a well-validated annotator makes errors, and errors that are correlated
with the variables under study bias the estimates and tests computed from the
labels. High agreement does not make the problem disappear: substituting
model labels directly into a regression has been shown to give substantial
bias and invalid confidence intervals even when the labels are 80 to 90\%
accurate \cite{egami2023using}. Methods now exist
that combine a small set of human labels with a large set of model labels to
produce estimates with valid confidence intervals: prediction-powered
inference \cite{angelopoulos2023prediction}, design-based supervised learning
for downstream regression \cite{egami2023using}, procedures that route
items to human coders according to the model's confidence
\cite{gligoric2025can}, and corrections for a judge's imperfect sensitivity
and specificity estimated from a human-labelled calibration set
\cite{lee2025how}. They suit HRI's typical situation of a few hundred to
a few thousand items of which a subset can be double-coded. At the least, the
main analysis should be repeated on the human-coded subset alone and any
divergence reported.

\subsection{R8 Judge}

\begin{practice}{bp:judge}{Validate a model judge per construct, keep it independent, and probe its biases}
Report judge--human agreement separately for each rubric dimension on
held-out material, use a judge from a different model family than the system
judged, test known biases, and keep humans as adjudicators of what the judge
flags.
\end{practice}

A Judge differs from an Annotator in what it scores: the outputs of another
model, usually against a rubric. The general evidence and its limits were
reviewed under \bpr{bp:validate}. The specific practices are these. Agreement
must be reported per dimension, because a judge that tracks human ratings of
fluency may not track ratings of safety. The judge should not share a model
family with the system it judges, given self-preference
\cite{panickssery2024llm}. Scoring should follow an explicit rubric with
reasoning elicited before the score \cite{liu2023geval}. Criteria should be
expected to change as outputs are inspected, a phenomenon documented as
criteria drift, which is a reason to iterate on the rubric with human raters
before freezing it \cite{shankar2024who}. Panels of diverse judges can reduce
the idiosyncrasy of any one \cite{verga2024replacing}. The form of the human
reference matters as well: where several ratings of an item are defensible,
forcing human raters to choose one biases the validation, and selecting
judges against such labels picked judges that performed markedly worse than
those selected against ratings that allowed a set of acceptable answers
\cite{guerdan2025validating}. For multimodal
material, vision-language judges have been found to track human preferences
in pairwise comparison better than in absolute scoring
\cite{chen2024mllm}, which favours comparative designs.

An evaluation of multi-agent judging for psychosocial safety is instructive
about trade-offs. Against human safety annotations, a pipeline in which a
second model reviewed the first model's judgement was the most accurate and
stable of four designs; a debate between adversarial models achieved
near-perfect recall at the price of over-flagging borderline cases; and
simple majority voting over repeated samples added little to a single
judgement \cite{luo2025dialogguard}. The appropriate design depends on the
cost of misses relative to false alarms, which in a Tier~III study is high
enough to favour recall, with human review of what is flagged. The authors'
conclusion, that judge scores are decision-support signals and not ground
truth, and their formative study with practitioners who used the tool as a
second opinion \cite{luo2026dialogguard}, describe the proper standing of a
Judge in HRI: a means of directing human attention, not of replacing it.

\subsection{R9 Simulator}

\begin{practice}{bp:simulator}{Use simulated people to stress systems, not to stand in for participants}
Restrict simulated users, partners, and raters to piloting, coverage testing,
and stress-testing; validate the simulator on the properties relied upon; and
do not offer simulated responses as evidence about human experience.
\end{practice}

Language models conditioned on personas can reproduce some aggregate patterns
of human survey responses and experimental effects
\cite{argyle2023out,aher2023using}, which has prompted proposals to use them
in place of participants \cite{dillion2023can}. The objections are by now
well developed. Simulated respondents show less variance than people and
their fidelity is unstable across prompts and over time
\cite{bisbee2024synthetic}; they misportray and flatten the identity groups
they are asked to represent \cite{wang2025large}; substituting them for
participants conflicts with the values of representation and inclusion that
motivate participatory research \cite{agnew2024illusion}; qualitative
researchers who tried them found that they foreclose participants' consent
and agency and lack contextual depth \cite{kapania2025simulacrum}; and the
case that they cannot replace participants has been made directly
\cite{harding2024ai}. For HRI there is a further and decisive point. The
discipline's subject is what happens when a person meets an embodied
machine, and a language model has never met one. The two attempts to
simulate HRI participants described under \bpr{bp:validate} bear this out:
simulated ratings tracked human ones for some measures and studies, failed
unpredictably for others, and inverted the dimension of comfort and
strangeness that is most bound up with physical presence
\cite{wachowiak2024are,lockwood2026when}.

Simulation nonetheless has legitimate uses. Before participants are
recruited, simulated users can exercise a dialogue system through many more
conversations than a pilot would, exposing failure modes, persona drift, and
unsafe responses. For this the simulator must itself be stable: multi-turn
simulation suffers from persona drift and from agents echoing each other, and
methods that address this have been evaluated with human validation
\cite{luo2026spasm}. What is found in simulation is a hypothesis about what
will happen with people, to be tested with people, and should be reported as
such.
